\documentclass{article}
\usepackage{subfig}
\usepackage{graphicx} % Required for inserting images
\usepackage{geometry}
\usepackage{hyperref}
\usepackage{amssymb}

\geometry{hmargin=2cm,vmargin=1.5cm}

\usepackage[english]{babel}
\usepackage[T1]{fontenc}

\title{Notes on energy conservation using the multilayer}
\author{Hugo Jacquet}
\date{21 July 2026}

\begin{document}

\maketitle

\section{Goal}

My goal is to study the interaction between a breaking wave field in different
environmental conditions. A first study will be on the interaction between
breaking, Langmuir cells and a stable stratification. For this, I need a
somewhate steady wave field, which is possible thanks to the wind forcing
described in ``Multilayer simulation of wind-forced breaking wave1
fields'' by Rui Yang and others.

\section{Setup}

Before looking at real sea state, I am working on a 2D wave to implement this
forcing. This setup goals are:
\begin{enumerate}

	\item verify that the multilayer model is able to reproduce the wave decay
	      predicted by Lamb (1932)

	\item implement the wind forcing

	\item verify that with the nudging term proposed by Rui Yang in his preprint I
	      can get a statically steady sea state

\end{enumerate}


My configuration is available at
\url{https://basilisk.fr/sandbox/hugoj/test_windinput/}. The first point is to
verify that energy decay following the exponential decay $E(t)/E_0=e^{-4 \nu k^2
	t}$. I initialize $\eta$ and the currents following the linear theory, then I
let the wave decay due to viscous dissipation. I use isotropic dissipation by
adding to the default vertical momentum diffusion horizontal diffusion with the
\verb|horizontal_diffusion| event.


My base case uses 10 periods for a box of 100 m at $R_e=400$ ($\nu=1/R_e=2e{-3}$
m2/s), with a depth of 5 m and a steepness $ak=0.01$. Numerics are using
defaults values (CFL=0.5, CFL\_H=1, theta\_H=0.5, theta=1.3, DT=HUGE). My small
script \verb|run_test_instrumented.sh| runs my base test case, the one from
Rui's preprint, and it also tests other values for different parameters.

Since I see some problem with energy conservation, I am looking for dissipative
processes in the numerics used. I use a version of the source code that allows
to count the number of time a gradient limiter is triggered, and the number of
time the limiter is actually applied. Output of this diagnostic is in the
\verb|log| file. I identified four limiters in the
multilayer code: i) A CFL limiter in \verb|hydro.h| that ensure positiveness of
the layer heights; ii) A general gradient limiter used in the BCG advection
scheme (minmod but with a counter, see the \verb|no_forcing_instrumented.c| file
where I redefine the gradient of scalars); iii) A limiter in the remap
polynomial reconstruction that avoid creating new extrema when remapping (in
\verb|remapc.h|); iv) A limiter that models the breaking by limiting the maximum
vertical velocity (in \verb|hydro.h| and \verb|nh.h|).


\section{Observations}

The base case shows a decay of energy that seems to match the theoretical decay.
However, when changing some parameters I observe that this is sometimes not the case
anymore, or even that the energy grows. \\

\subsection{Base case results}

Figure \ref{fig:base} shows the normalized energy decay for my base case.
Numerical and analytical results are qualitatively close, energy decay (as it
should). Explicit viscosity seems to be the driver for the dissipation of the
monochromatic wave.


\begin{figure}
	\centering
	\includegraphics[width=.5\linewidth]{results_instrumented/base/energy.png}
	\caption{Base case results}
	\label{fig:base}
\end{figure}

\subsection{Rui's case is growing in energy despite no forcing}

Figure \ref{fig:Ruis} clearly shows that something is wrong with my code in this
configuration. Both kinetic and potential energy increase, which should not
happen in this no-forcing configuration.

\begin{figure}
	\center
	\includegraphics[width=0.5\textwidth]{results_instrumented/Ruis_thetah0.5/energy.png}
	\caption{Normalized energy evolution for parameters from Rui's preprint}
	\label{fig:Ruis}
\end{figure}

\subsection{Base case seems converged, but it's not}

In the base case (and this is the default behavior in Basilisk), the maximum
time-step is set to \verb|HUGE|, which is just a big number, and the actual
time-step used in the solver is defined by the CFL condition(s). There are two
CFL in the multilayer, one related to the barotropic part of the solver
(advection of $\eta$) and one related to the non-hydrostatic part (advection of
currents $u$ and $v$). A discussion about the influence of the CFL is below. The
user can restrict the value of the maximum time-step by prescribing \verb|DT|.
Figure \ref{fig:changeDT} shows the influence of this maximum time-step on the
energy evolution. When \verb|DT| is automatically chosen (brown line), the
energy decay is close to the theoretical value and to the value \verb|DT| = 0.1
s. When inspecting the actual, variable time-step it is close to 0.1 s, so we
expect the purple and brown line to be close to each other. However, when
decreasing the maximum time-step I observe that dissipation decreases. Reducing
numerical dissipation when increasing (temporal) resolution is expected, but my
base case seems to be close to the theoretical expression, and so I expected it
to be already converged. I am now wondering if this base case is really
converged or is it just ``luck'' that I get an energy evolution close to the
analytic solution...

In the same reasoning, increasing resolution through $N$ should thus not change
the result if the base case is converged. Figure \ref{fig:changeN} shows that
this is not the case as when increasing the spatial resolution the dissipation
decreases to a point where the highest resolution simulation, which I expect to
be the most resolved, is not as dissipative as the theoretical expression. I
expect the low resolution simulations to be over dissipative due to the combined
effect of the numerical and explicit diffusions, but the fact that when explicit
diffusion is the main dissipation mechanism the energy decay doesn't match the
prediction is beyond my understanding. When increasing the temporal resolution,
we get the first observation.

\begin{figure}
	\center
	\includegraphics[width=0.5\textwidth]{results_instrumented/change_of_DT.png}
	\caption{Normalized energy evolution for the base set of parameters but the maximum
		time-step is limited by hand and decreased}
	\label{fig:changeDT}
\end{figure}


\begin{figure}
	\center
	\includegraphics[width=0.5\textwidth]{results_instrumented/change_of_N.png}
	\caption{Normalized energy decay with varying horizontal resolution}
	\label{fig:changeN}
\end{figure}

\subsection{A well resolved case}

I need a test case where I know the solution is resolved. If this case is not
well reproduced by the code, then I need to first fix this problem. I thus build
a case that should be resolved with 100 points per wavelength, 100 points per
period and a Re high enough to get a scale separation between viscous
dissipation and wave dynamic.



\subsection{Change of regime no longer matches the analytical law}

In an effort to bridge between the set of parameters from Rui's preprint and my
base case, I investigated the influence of the value of the explicit viscosity
on the solution. For this test, I use a linear wave with 1 wavelength (compared
to the test before with 10 wavelengths).




\end{document}

